
\documentclass[../../main.tex]{subfiles}
\begin{document}

% 年度の大きな表示
\begin{center}
  {\fontsize{48pt}{58pt}\selectfont \textbf{1983}\normalsize\textbf{年度}}
\end{center}

\vspace{10mm}

% 本文

\vspace{15mm}

% 出題テーマと難易度の表
\noindent\textbf{出題テーマと難易度}

\vspace{5mm}

\begin{center}
  \shadowbox{%
    \begin{tabular}{|c|p{8cm}|c|}
      \hline
      \textbf{問題番号} & \textbf{テーマ} & \textbf{難易度} \\
      \hline
      第1問 & 整数論と極限 & \\
      \hline
      第2問 & 曲線群の通過領域 & \\
      \hline
      第3問 & 行列と漸化式 & \\
      \hline
      第4問 & 三角関数の最大・最小 & \\
      \hline
      第5問 & 双曲線と面積の最小値 & \\
      \hline
    \end{tabular}%
  }
\end{center}

\vspace{15mm}

% 解答
\noindent\textbf{解答}

\vspace{5mm}

\begin{center}
  \shadowbox{%
    \renewcommand{\arraystretch}{1.6}
    \begin{tabular}{|c|c|p{8cm}|}
      \hline
      \textbf{問題番号} & \textbf{小問} & \textbf{答え} \\
      \hline
      \multirow{2}{*}{第1問} & (1) & 証明問題 \\
      \cline{2-3}
                             & (2) & $\dfrac{15}{2}$ \\
      \hline
      第2問 & - & 図示（$0<x<2$部分は半円/放物線の上方，$x\le0,2\le x$部分は放物線$y=-(x-1)^2+\dfrac{5}{4}$の下方の領域） \\
      \hline
      第3問 & - &  \\
      \hline
      \multirow{2}{*}{第4問} & (1) & 最大値 $\dfrac{3\sqrt{3}}{2}$ \\
      \cline{2-3}
                             & (2) & $-2<F\le\dfrac{3\sqrt{3}}{2}$ \\
      \hline
      第5問 & - & 最小値 $8+6\log3$（$t=3$） \\
      \hline
    \end{tabular}%
    \renewcommand{\arraystretch}{1.0}
  }
\end{center}

\vspace{15mm}

% 解答時間
\noindent\textbf{試験形態}

\vspace{5mm}

\begin{center}
  \shadowbox{%
    \begin{tabular}{p{8cm}}
      （要確認）
    \end{tabular}%
  }
\end{center}

\end{document}
